% preamble-exam-zh.tex — 共享 preamble
% 用于所有 exam-zh 模板
%
% 样式策略：
%   屏幕 → 语义彩色（蓝=关键想法，红=易错，绿=路线，灰=提示/教师备注）
%   打印 → 自动降级为黑白（通过 print mode 或 monochrome 选项）

\documentclass{exam-zh}

% ---------- 基础配置 ----------
\examsetup{
  page/size = a4paper,
  page/show-head = false,
  page/show-foot = false,
  sealline/show = false,
  font = times,
  math-font = xits,
}

\usepackage{tcolorbox}
\usepackage{needspace}
\tcbuselibrary{skins,breakable}

% ---------- 打印/屏幕颜色切换 ----------
% 用法：\documentclass[monochrome]{exam-zh} 即可切换为纯黑白
% 注意：不要使用 \if@xxx 放在 makeatother 外部；这里改成普通条件名。
\newif\ifeduprintcolor
\eduprintcolortrue

\makeatletter
\@ifclasswith{exam-zh}{monochrome}{\eduprintcolorfalse}{}
\makeatother

% ---------- 语义颜色定义 ----------
% 屏幕：有色彩区分；打印：降级为灰度
\ifeduprintcolor
  % --- 屏幕彩色模式 ---
  \definecolor{edu-blue}{RGB}{47, 82, 122}       % 关键想法
  \definecolor{edu-blue-bg}{RGB}{243, 246, 252}
  \definecolor{edu-red}{RGB}{180, 40, 40}         % 易错提醒
  \definecolor{edu-red-bg}{RGB}{255, 245, 240}
  \definecolor{edu-green}{RGB}{34, 100, 60}       % 路线图
  \definecolor{edu-green-bg}{RGB}{242, 250, 245}
  \definecolor{edu-orange}{RGB}{160, 90, 0}       % 训练目标
  \definecolor{edu-orange-bg}{RGB}{255, 248, 235}
  \definecolor{edu-gray}{RGB}{100, 100, 100}      % 提示/教师备注
  \definecolor{edu-gray-bg}{RGB}{248, 248, 248}
  \definecolor{edu-purple}{RGB}{90, 50, 120}      % 思考题
  \definecolor{edu-purple-bg}{RGB}{248, 244, 252}
\else
  % --- 打印黑白模式 ---
  \definecolor{edu-blue}{gray}{0.15}
  \definecolor{edu-blue-bg}{gray}{0.96}
  \definecolor{edu-red}{gray}{0.25}
  \definecolor{edu-red-bg}{gray}{0.97}
  \definecolor{edu-green}{gray}{0.20}
  \definecolor{edu-green-bg}{gray}{0.97}
  \definecolor{edu-orange}{gray}{0.25}
  \definecolor{edu-orange-bg}{gray}{0.98}
  \definecolor{edu-gray}{gray}{0.40}
  \definecolor{edu-gray-bg}{gray}{0.97}
  \definecolor{edu-purple}{gray}{0.30}
  \definecolor{edu-purple-bg}{gray}{0.98}
\fi

% ---------- 自定义教学环境 ----------
% 共性：enhanced + breakable（跨页不断裂）

% 原题卡片 — 强边框，突出显示
\newtcolorbox{problemcard}{
  enhanced, breakable,
  colback=edu-blue-bg,
  colframe=edu-blue,
  boxrule=1.2pt,
  arc=0pt,
  left=3mm, right=3mm, top=3mm, bottom=3mm,
}

% 解题路线图 — 左边线 + 浅底
\newtcolorbox{routemap}{
  enhanced, breakable,
  colback=edu-green-bg,
  colframe=edu-green!30,
  boxrule=0pt,
  borderline west={2.5pt}{0pt}{edu-green},
  left=3mm, right=3mm, top=3mm, bottom=3mm,
}

% 关键想法 — 蓝色左边线
\newtcolorbox{keyideabox}{
  enhanced, breakable,
  colback=edu-blue-bg,
  colframe=edu-blue!30,
  boxrule=0pt,
  borderline west={2.5pt}{0pt}{edu-blue},
  left=3mm, right=3mm, top=3mm, bottom=3mm,
}

% 易错提醒 — 红色左边线
\newtcolorbox{mistakebox}{
  enhanced, breakable,
  colback=edu-red-bg,
  colframe=edu-red!20,
  boxrule=0pt,
  borderline west={3pt}{0pt}{edu-red},
  left=3mm, right=3mm, top=2.5mm, bottom=2.5mm,
}

% 提示 — 灰色虚线左边线
\newtcolorbox{hintbox}{
  enhanced, breakable,
  colback=edu-gray-bg,
  colframe=edu-gray!20,
  boxrule=0pt,
  borderline west={2pt}{0pt}{edu-gray, dashed},
  left=3mm, right=3mm, top=2mm, bottom=2mm,
}

% 思考题 — 紫色虚线左边线
\newtcolorbox{thinkbox}{
  enhanced, breakable,
  colback=edu-purple-bg,
  colframe=edu-purple!20,
  boxrule=0pt,
  borderline west={2pt}{0pt}{edu-purple, dashed},
  left=2.5mm, right=2.5mm, top=2.5mm, bottom=2.5mm,
}

% 教师备注 — 灰色左边线，小字
\newtcolorbox{teachernote}{
  enhanced, breakable,
  colback=edu-gray-bg,
  colframe=edu-gray!15,
  boxrule=0pt,
  borderline west={2pt}{0pt}{edu-gray!60},
  left=2.5mm, right=2.5mm, top=2.5mm, bottom=2.5mm,
  fontupper=\small,
  colupper=black!60,
}

% 训练目标 — 橙色左边线，小字
\newtcolorbox{traininggoal}{
  enhanced, breakable,
  colback=edu-orange-bg,
  colframe=edu-orange!15,
  boxrule=0pt,
  borderline west={1.5pt}{0pt}{edu-orange},
  left=2mm, right=2mm, top=1mm, bottom=1mm,
  fontupper=\small,
}

% 答题区 — 无边框，纯留白
\newcommand{\answerarea}[1][40mm]{%
  \vspace{2mm}%
  \begin{tcolorbox}[
    enhanced, breakable,
    colback=white,
    colframe=white,
    boxrule=0pt,
    height=#1,
    valign=top,
    bottom=0mm,
    after skip=0mm,
  ]
  \end{tcolorbox}%
}

% ---------- 字体设置 ----------
% exam-zh (ctexbook) already sets CJK fonts via ctex.
% Only override if specific fonts are needed.
% macOS: Songti SC, STHeiti, STFangsong
% Linux: SimSun, SimHei, FangSong

% ---------- 页面设置 ----------
% exam-zh (ctexbook) already loads geometry.
% Override margins if needed:
% \geometry{top=16mm, bottom=16mm, left=14mm, right=14mm}

\examsetup{
  question/show-answer = false,
  fillin/show-answer = false,
  solution/show-solution = hide,
}

% ---------- 教师版解答样式 ----------
\newtcolorbox{solutionblock}{
  enhanced, breakable,
  colback=edu-blue-bg,
  colframe=edu-blue!25,
  boxrule=0pt,
  borderline west={2pt}{0pt}{edu-blue!50},
  arc=0pt,
  left=3mm, right=3mm, top=2mm, bottom=2mm,
  before skip=2mm,
  after skip=3mm,
  fontupper=\small,
}

% 解答步骤编号
\newcounter{solstep}
\newcommand{\solstep}[2]{%
  \stepcounter{solstep}%
  \par\medskip
  \noindent\hangindent=6mm
  {\color{edu-blue}\bfseries\textcircled{\scriptsize\thesolstep}\enspace #1}\enspace #2\par
}

\begin{document}

\section*{已知两点求一次函数解析式 · 专题练习}

\section*{一、填空题（每题 5 分，共 30 分）}


\begin{question}[points=5]
  一次函数的图像经过 $(1, 5)$ 和 $(-1, 1)$ 两点，则这个一次函数的解析式为 \fillin。
  \fillin
\end{question}



\begin{question}[points=5]
  一次函数 $y = kx + b$ 与 $y = 3x - 1$ 平行，且经过点 $(1, 2)$，则 $b = $ \fillin。
  \fillin
\end{question}



\begin{question}[points=5]
  将直线 $y = -3x + 5$ 向上平移 $4$ 个单位后，所得直线的解析式为 \fillin。
  \fillin
\end{question}



\begin{question}[points=5]
  直线 $y = 3x - 6$ 关于 $y$ 轴对称的直线解析式为 \fillin。
  \fillin
\end{question}



\begin{question}[points=5]
  已知一次函数 $y = kx + b$ 与 $y = -x + 4$ 没有交点，且在 $y$ 轴上的截距为 $7$，则解析式为 \fillin。
  \fillin
\end{question}



\begin{question}[points=5]
  已知 $f(x)$ 为一次函数，$f(x + 1) + f(x - 1) = 4x + 6$，则 $f(x) = $ \fillin。
  \fillin
\end{question}


\section*{二、选择题（每题 5 分，共 30 分）}


\begin{question}[points=5]
  一次函数 $y = kx + b$ 的图像经过 $(0, -3)$ 和 $(2, 1)$，则 $k$ 的值为
  \begin{choices}
    \item $-2$
    \item $2$
    \item $-1$
    \item $1$
  \end{choices}
\end{question}



\begin{question}[points=5]
  下列直线中，与 $y = 2x + 3$ 平行的是
  \begin{choices}
    \item $y = -2x + 1$
    \item $y = 2x - 5$
    \item $y = \frac{1}{2}x + 3$
    \item $y = x + 3$
  \end{choices}
\end{question}



\begin{question}[points=5]
  将直线 $y = x + 2$ 向左平移 $3$ 个单位，再向下平移 $1$ 个单位，所得直线的解析式为
  \begin{choices}
    \item $y = x - 2$
    \item $y = x + 4$
    \item $y = x + 1$
    \item $y = x - 1$
  \end{choices}
\end{question}



\begin{question}[points=5]
  直线 $y = -x + 4$ 关于 $x$ 轴对称的直线解析式为
  \begin{choices}
    \item $y = x + 4$
    \item $y = x - 4$
    \item $y = -x - 4$
    \item $y = -x + 4$
  \end{choices}
\end{question}



\begin{question}[points=5]
  一次函数 $y = kx + b$ 满足 $f(0) = 3$，$f(2) = -1$，则 $f(-1) = $
  \begin{choices}
    \item $4$
    \item $5$
    \item $6$
    \item $7$
  \end{choices}
\end{question}



\begin{question}[points=5]
  直线 $y = kx + b$ 与 $y = 2x + 1$ 无交点，且经过点 $(-1, 5)$，则 $b$ 的值为
  \begin{choices}
    \item $7$
    \item $3$
    \item $-7$
    \item $5$
  \end{choices}
\end{question}


\section*{三、解答题（每题 10 分，共 40 分）}

\needspace{8\baselineskip}

\begin{problem}[points=10]
  将直线 $y = -x + 3$ 向右平移 $2$ 个单位，再向上平移 $1$ 个单位，得到直线 $l_1$。再将 $l_1$ 关于 $y$ 轴对称，得到直线 $l_2$。求 $l_2$ 的解析式。
  \answerarea[80mm]
\end{problem}


\needspace{8\baselineskip}

\begin{problem}[points=10]
  直线 $y = -2x + 6$ 与 $x$ 轴交于点 $A$，与 $y$ 轴交于点 $B$。直线 $l$ 过原点 $O$ 且与线段 $AB$ 交于点 $C$，若 $\triangle AOC$ 与 $\triangle BOC$ 的面积之比为 $1 : 2$，求直线 $l$ 的解析式。
  \answerarea[100mm]
\end{problem}


\needspace{8\baselineskip}

\begin{problem}[points=10]
  已知一次函数 $y = kx + b$ 与 $y = -3x + 2$ 平行，且与直线 $y = x - 4$ 交于 $y$ 轴上的同一点，求这个一次函数的解析式。
  \answerarea[60mm]
\end{problem}


\needspace{8\baselineskip}

\begin{problem}[points=10]
  已知一次函数 $y = kx + b$（$k \neq 0$），当 $-1 \leq x \leq 2$ 时，$y$ 的最大值与最小值之差为 $9$。当 $x = 0$ 时 $y = 3$。求这个一次函数的解析式。
  \answerarea[80mm]
\end{problem}



\end{document}
